\documentclass[12pt,a4paper]{report}

% ─── Packages ───────────────────────────────────────────────────────────────
\usepackage[a4paper, margin=1in]{geometry}
\usepackage{graphicx}
\usepackage{amsmath, amssymb}
\usepackage{booktabs}
\usepackage{array}
\usepackage{longtable}
\usepackage{multirow}
\usepackage{caption}
\usepackage{subcaption}
\usepackage{float}
\usepackage{fancyhdr}
\usepackage{titlesec}
\usepackage{tocloft}
\usepackage{hyperref}
\usepackage{xcolor}
\usepackage{listings}
\usepackage{mdframed}
\usepackage{enumitem}
\usepackage{tikz}
\usepackage{pgfplots}
\pgfplotsset{compat=1.18}
\usepackage{colortbl}
\usepackage{tabularx}
\usepackage{makecell}
\usepackage{setspace}
\usepackage{parskip}
\usepackage{microtype}
\usepackage{fontenc}
\usepackage{inputenc}

% ─── Hyperref Setup ──────────────────────────────────────────────────────────
\hypersetup{
    colorlinks=true,
    linkcolor=blue!70!black,
    urlcolor=cyan!70!black,
    citecolor=green!60!black,
    pdftitle={Design and Implementation of Universal Shift Register using Verilog},
    pdfauthor={},
    pdfsubject={Digital Systems Design},
}

% ─── Color Palette ───────────────────────────────────────────────────────────
\definecolor{codegreen}{rgb}{0.13,0.55,0.13}
\definecolor{codegray}{rgb}{0.5,0.5,0.5}
\definecolor{codepurple}{rgb}{0.58,0.0,0.83}
\definecolor{codeblue}{rgb}{0.13,0.29,0.53}
\definecolor{codered}{rgb}{0.75,0.10,0.05}
\definecolor{backcolour}{rgb}{0.97,0.97,0.97}
\definecolor{keywordcolor}{rgb}{0.0,0.3,0.7}
\definecolor{titleblue}{RGB}{0,70,127}
\definecolor{sectionblue}{RGB}{0,90,160}
\definecolor{lightgray}{RGB}{240,240,240}
\definecolor{darkgray}{RGB}{80,80,80}
\definecolor{tableheader}{RGB}{0,70,127}

% ─── Listings (Verilog) ──────────────────────────────────────────────────────
\lstdefinelanguage{Verilog}{
    morekeywords=[1]{
        module, endmodule, input, output, reg, wire, always, begin, end,
        if, else, case, endcase, posedge, negedge, assign, parameter,
        localparam, initial, $dumpfile, $dumpvars, $finish, $monitor,
        $display, $time, integer, for, default, 'b0, 'b1
    },
    morekeywords=[2]{
        and, or, not, nand, nor, xor, xnor, buf
    },
    sensitive=true,
    morecomment=[l]{//},
    morecomment=[s]{/*}{*/},
    morestring=[b]",
}

\lstdefinestyle{verilogstyle}{
    language=Verilog,
    backgroundcolor=\color{backcolour},
    commentstyle=\color{codegreen}\itshape,
    keywordstyle=[1]\color{keywordcolor}\bfseries,
    keywordstyle=[2]\color{codepurple}\bfseries,
    numberstyle=\tiny\color{codegray},
    stringstyle=\color{codered},
    basicstyle=\ttfamily\footnotesize,
    breakatwhitespace=false,
    breaklines=true,
    captionpos=b,
    keepspaces=true,
    numbers=left,
    numbersep=10pt,
    showspaces=false,
    showstringspaces=false,
    showtabs=false,
    tabsize=4,
    frame=single,
    framerule=0.8pt,
    rulecolor=\color{codeblue},
    xleftmargin=15pt,
    framexleftmargin=15pt,
}

\lstset{style=verilogstyle}

% ─── Header / Footer ─────────────────────────────────────────────────────────
\pagestyle{fancy}
\fancyhf{}
\fancyhead[L]{\small\textcolor{darkgray}{Universal Shift Register}}
\fancyhead[R]{\small\textcolor{darkgray}{Verilog Implementation}}
\fancyfoot[C]{\thepage}
\fancyfoot[R]{\small\textcolor{darkgray}{\today}}
\renewcommand{\headrulewidth}{0.5pt}
\renewcommand{\footrulewidth}{0.3pt}

% ─── Section Formatting ──────────────────────────────────────────────────────
\titleformat{\chapter}[display]
  {\normalfont\huge\bfseries\color{titleblue}}
  {\chaptertitlename\ \thechapter}{18pt}{\Huge}

\titleformat{\section}
  {\normalfont\Large\bfseries\color{sectionblue}}
  {\thesection}{1em}{}

\titleformat{\subsection}
  {\normalfont\large\bfseries\color{darkgray}}
  {\thesubsection}{1em}{}

\titleformat{\subsubsection}
  {\normalfont\normalsize\bfseries}
  {\thesubsubsection}{1em}{}

% ─── Table of Contents Formatting ────────────────────────────────────────────
\renewcommand{\cftchapfont}{\bfseries\color{titleblue}}
\renewcommand{\cftsecfont}{\color{sectionblue}}
\setlength{\cftbeforechapskip}{6pt}

% ─── Spacing ─────────────────────────────────────────────────────────────────
\onehalfspacing

% ─── Custom Commands ─────────────────────────────────────────────────────────
\newcommand{\signal}[1]{\texttt{\textcolor{codeblue}{#1}}}
\newcommand{\keyword}[1]{\textbf{\textcolor{keywordcolor}{#1}}}
\newcommand{\reg}[1]{\texttt{#1}}
\newcommand{\highlight}[1]{\colorbox{lightgray}{\texttt{#1}}}

% ─── mdframed box for notes ──────────────────────────────────────────────────
\newmdenv[
  topline=true, bottomline=true, rightline=false, leftline=true,
  linewidth=3pt, linecolor=sectionblue,
  backgroundcolor=blue!4,
  innerleftmargin=10pt, innerrightmargin=10pt,
  innertopmargin=8pt, innerbottommargin=8pt,
]{notebox}

\newmdenv[
  topline=true, bottomline=true, rightline=false, leftline=true,
  linewidth=3pt, linecolor=codered,
  backgroundcolor=red!4,
  innerleftmargin=10pt, innerrightmargin=10pt,
  innertopmargin=8pt, innerbottommargin=8pt,
]{warnbox}

% ═══════════════════════════════════════════════════════════════════════════════
%                              DOCUMENT START
% ═══════════════════════════════════════════════════════════════════════════════
\begin{document}

% ─── Title Page ──────────────────────────────────────────────────────────────
\begin{titlepage}
    \centering
    \vspace*{1cm}

    {\Huge\bfseries\color{titleblue}
        Design and Implementation of\\[6pt]
        Universal Shift Register\\[6pt]
        Using Verilog HDL\par}

    \vspace{0.5cm}
    \rule{\linewidth}{1.5pt}
    \vspace{0.3cm}

    {\large\color{darkgray}
        A Comprehensive Technical Report on the Architecture,\\
        Working Principle, RTL Design, and Simulation of a\\
        4-Bit Universal Shift Register\par}

    \vspace{0.5cm}
    \rule{\linewidth}{0.5pt}

    \vspace{2cm}

    \begin{mdframed}[linecolor=sectionblue, linewidth=1pt, roundcorner=5pt]
    \centering
    \begin{tabular}{ll}
        \textbf{Subject}        & : Digital Systems Design / VLSI \& HDL Lab \\[4pt]
        \textbf{Topic}          & : Universal Shift Register (4-bit) \\[4pt]
        \textbf{HDL Language}   & : Verilog (IEEE 1364-2001 Standard) \\[4pt]
        \textbf{Simulator}      & : ModelSim / Icarus Verilog (iverilog) \\[4pt]
        \textbf{Date}           & : \today \\
    \end{tabular}
    \end{mdframed}

    \vfill

    \includegraphics[width=0.18\textwidth]{logo.png} % Replace with your institution logo
    \vspace{0.5cm}

    {\large\bfseries Department of Electronics \& Communication Engineering\\[4pt]
    \textit{[Institution Name]}\par}

    \vspace{0.5cm}
    {\small\color{darkgray} All simulation results and waveforms were generated using Icarus Verilog and GTKWave.}
\end{titlepage}

% ─── Abstract ────────────────────────────────────────────────────────────────
\chapter*{Abstract}
\addcontentsline{toc}{chapter}{Abstract}

A \textbf{Universal Shift Register (USR)} is one of the most versatile and foundational building blocks in digital design. Unlike a conventional shift register that is limited to a single mode of operation, a universal shift register integrates four distinct operating modes — \textit{parallel load}, \textit{shift right}, \textit{shift left}, and \textit{hold} — into a single, unified circuit controlled by two select lines. This multi-functional nature makes it indispensable in data communication systems, arithmetic logic units, serial-to-parallel and parallel-to-serial converters, and a wide range of sequential digital systems.

This report presents a thorough study of the architecture and internal workings of a 4-bit universal shift register, followed by its complete implementation using the Verilog Hardware Description Language (HDL). The design is described at the Register Transfer Level (RTL), making it synthesisable for FPGAs and ASICs. A complete Verilog design module along with a self-checking testbench that exercises all four modes of operation is provided. The behaviour of the circuit is examined through detailed timing diagram analysis, and each section of the code is explained line-by-line.

The key contributions of this report are:
\begin{itemize}[leftmargin=2em]
    \item Detailed explanation of the internal D flip-flop and multiplexer architecture.
    \item Full truth/state table covering all four operating modes.
    \item RTL-level Verilog design code with inline commentary.
    \item A comprehensive testbench that applies all modes with diverse input vectors.
    \item Annotated timing diagram walkthrough demonstrating functional correctness.
\end{itemize}

\newpage

% ─── Table of Contents ───────────────────────────────────────────────────────
\tableofcontents
\newpage
\listoffigures
\listoftables
\newpage

% ═══════════════════════════════════════════════════════════════════════════════
%  CHAPTER 1: Introduction
% ═══════════════════════════════════════════════════════════════════════════════
\chapter{Introduction}

\section{Background and Motivation}

In digital electronics, the ability to store, transfer, and manipulate binary data in a controlled and sequential manner is fundamental to the operation of virtually every modern computing system. Registers are the core storage elements that serve this purpose at the heart of processor datapaths, communication interfaces, and control units.

A \textbf{shift register} is a cascade of flip-flops that shares a common clock signal; each flip-flop transfers its stored bit to the next stage at every rising (or falling) clock edge. While simple shift registers provide unidirectional data movement, practical systems routinely demand the ability to move data in both directions and to load or read multiple bits in parallel simultaneously.

The \textbf{Universal Shift Register} (USR) was conceived to meet all these requirements in a single device. The 74HC194 and 74LS194 are classic industry-standard universal shift register ICs that have been widely used since the 1970s. By incorporating a 4-to-1 multiplexer in front of each flip-flop and exposing two mode-select lines, the USR can be programmed on a clock-cycle-by-clock-cycle basis to perform any of its four fundamental operations.

\section{Scope of the Report}

This report focuses on the design of a \textbf{4-bit universal shift register}. A 4-bit design is the canonical reference width — it is wide enough to demonstrate all inter-bit interactions, yet compact enough that truth tables and timing diagrams remain readable. The design is fully parameterisable and can be extended to 8, 16, or 32 bits by changing a single parameter.

The implementation language is \textbf{Verilog HDL} (IEEE Std 1364-2001), the most widely adopted hardware description language in industry and academia for RTL design. The simulation is performed using Icarus Verilog, and waveform viewing is done with GTKWave, both of which are freely available open-source tools.

\section{Objectives}

The primary objectives of this report are as follows:

\begin{enumerate}[leftmargin=2.5em, label=\arabic*.]
    \item To understand the architecture of a universal shift register from first principles.
    \item To derive and document the complete mode-select truth table and state behaviour.
    \item To implement the design in Verilog HDL at the Register Transfer Level.
    \item To verify functional correctness using a comprehensive testbench that covers all four modes.
    \item To analyse the simulation waveform and correlate it with the expected logical behaviour.
    \item To identify practical applications where a universal shift register is deployed.
\end{enumerate}

\section{Historical Context and Industrial Relevance}

The universal shift register has its roots in the TTL logic family of the early 1970s. Texas Instruments introduced the SN74194 as a 4-bit bidirectional universal shift register with parallel access, which became one of the most popular off-the-shelf sequential logic ICs. Its descendants — the 74HC194, 74HCT194, and 74F194 — are still in production today and are widely used in embedded system prototyping.

In modern FPGA and ASIC design flows, the shift register is typically not instantiated as an off-the-shelf IC but is rather described at the RTL level in Verilog or VHDL and synthesised by a tool such as Xilinx Vivado, Intel Quartus, or Synopsys Design Compiler. The Verilog description presented in this report follows the RTL coding style that is directly compatible with these industry synthesis tools.

\begin{notebox}
\textbf{Note:} The 4-bit width used throughout this report is chosen for pedagogical clarity. All design constructs use the \texttt{N} parameter, and widening to 8 or 16 bits requires only a single-line change.
\end{notebox}

% ═══════════════════════════════════════════════════════════════════════════════
%  CHAPTER 2: Components Needed
% ═══════════════════════════════════════════════════════════════════════════════
\chapter{Components Needed}

\section{Overview}

The universal shift register is constructed entirely from two categories of digital building blocks: \textbf{edge-triggered D flip-flops} for sequential storage, and \textbf{4-to-1 multiplexers} for combinational mode selection. Understanding each component's function individually is essential before grasping how they interact in the assembled circuit.

\section{D Flip-Flop (D-FF)}

\subsection{Definition and Symbol}

A \textbf{D flip-flop} (Data or Delay flip-flop) is a synchronous bistable element with a single data input \signal{D}, a clock input \signal{CLK}, and two complementary outputs \signal{Q} and \signal{\={Q}}. On the active clock edge, the flip-flop captures the logic value present at \signal{D} and holds it at the output \signal{Q} until the next active edge.

\subsection{Characteristic Equation}

The characteristic equation of the D flip-flop is:
\begin{equation}
    Q^{n+1} = D
    \label{eq:dff}
\end{equation}

This means the next state of the output equals the current input unconditionally — the simplest possible flip-flop behaviour, which makes it ideal for use in shift registers.

\subsection{Truth Table of D Flip-Flop}

\begin{table}[H]
    \centering
    \caption{Truth Table of D Flip-Flop (Positive Edge-Triggered)}
    \begin{tabular}{
        >{\centering\arraybackslash}m{2.5cm}
        >{\centering\arraybackslash}m{2.5cm}
        >{\centering\arraybackslash}m{3cm}
        >{\centering\arraybackslash}m{4cm}}
        \toprule
        \rowcolor{tableheader}
        \textcolor{white}{\textbf{CLK Edge}} &
        \textcolor{white}{\textbf{D}} &
        \textcolor{white}{\textbf{Q\textsuperscript{n+1}}} &
        \textcolor{white}{\textbf{Description}} \\
        \midrule
        $\uparrow$ (Rising) & 0 & 0 & Output resets to 0 \\
        $\uparrow$ (Rising) & 1 & 1 & Output sets to 1 \\
        No edge & $\times$ & $Q^n$ & Output holds \\
        \bottomrule
    \end{tabular}
    \label{tab:dff}
\end{table}

\subsection{Role in Shift Register}

In a universal shift register, \textbf{one D flip-flop is required per bit}. A 4-bit USR therefore contains exactly \textbf{four D flip-flops}, labelled FF3, FF2, FF1, and FF0, where FF3 holds the most significant bit (MSB) and FF0 holds the least significant bit (LSB).

\section{4-to-1 Multiplexer (MUX)}

\subsection{Definition}

A \textbf{4-to-1 multiplexer} is a combinational circuit with four data inputs (\signal{I0}–\signal{I3}), two select lines (\signal{S1}, \signal{S0}), and one output (\signal{Y}). It routes exactly one of the four inputs to the output based on the binary value encoded by the select lines.

\subsection{Boolean Expression}

\begin{equation}
    Y = \overline{S_1}\,\overline{S_0}\,I_0 \;+\; \overline{S_1}\,S_0\,I_1 \;+\; S_1\,\overline{S_0}\,I_2 \;+\; S_1\,S_0\,I_3
    \label{eq:mux}
\end{equation}

\subsection{Truth Table of 4-to-1 MUX}

\begin{table}[H]
    \centering
    \caption{Truth Table of 4-to-1 Multiplexer}
    \begin{tabular}{
        >{\centering\arraybackslash}m{2cm}
        >{\centering\arraybackslash}m{2cm}
        >{\centering\arraybackslash}m{4cm}}
        \toprule
        \rowcolor{tableheader}
        \textcolor{white}{\textbf{S1}} &
        \textcolor{white}{\textbf{S0}} &
        \textcolor{white}{\textbf{Output Y}} \\
        \midrule
        0 & 0 & I0 \\
        0 & 1 & I1 \\
        1 & 0 & I2 \\
        1 & 1 & I3 \\
        \bottomrule
    \end{tabular}
    \label{tab:mux}
\end{table}

\subsection{Role in Shift Register}

In the universal shift register, \textbf{one 4-to-1 MUX is placed in front of each D flip-flop}. The output of each MUX feeds the \signal{D} input of its associated flip-flop. The four inputs of each MUX are connected as follows:

\begin{itemize}
    \item \textbf{I0} $\rightarrow$ Current bit output $Q_i$ (Hold mode — the bit feeds back into itself)
    \item \textbf{I1} $\rightarrow$ Output of the bit to the left, $Q_{i+1}$ (Shift Right — MSB feeds LSB)
    \item \textbf{I2} $\rightarrow$ Output of the bit to the right, $Q_{i-1}$ (Shift Left — LSB feeds MSB)
    \item \textbf{I3} $\rightarrow$ Parallel data input $P_i$ (Parallel Load)
\end{itemize}

\section{Additional Components}

\subsection{Clock Source}

A periodic square wave clock signal is needed to synchronise all flip-flops. In simulation, this is modelled using a \texttt{forever} loop with a half-period delay. In hardware, a dedicated crystal oscillator or PLL output provides this.

\subsection{Reset (Clear) Logic}

An active-low or active-high asynchronous or synchronous reset line is used to initialise all flip-flop outputs to logic zero at power-up or at any time during operation.

\subsection{Serial Input Lines}

Two serial data entry pins are required:
\begin{itemize}
    \item \signal{MSB\_IN} (also called \signal{SIN\_R}): The serial bit shifted into the MSB position (FF3) during Shift Right mode.
    \item \signal{LSB\_IN} (also called \signal{SIN\_L}): The serial bit shifted into the LSB position (FF0) during Shift Left mode.
\end{itemize}

\subsection{Summary of Components for 4-bit USR}

\begin{table}[H]
    \centering
    \caption{Component Summary for 4-bit Universal Shift Register}
    \begin{tabular}{lll}
        \toprule
        \rowcolor{tableheader}
        \textcolor{white}{\textbf{Component}} &
        \textcolor{white}{\textbf{Quantity}} &
        \textcolor{white}{\textbf{Purpose}} \\
        \midrule
        D Flip-Flop (positive edge) & 4 & Bit storage (FF0 – FF3) \\
        4-to-1 Multiplexer & 4 & Mode selection per bit \\
        Clock Buffer & 1 & Distributes clock to all FFs \\
        Serial Input (Right) \signal{MSB\_IN} & 1 & Data entry in shift-right mode \\
        Serial Input (Left) \signal{LSB\_IN} & 1 & Data entry in shift-left mode \\
        Parallel Data Bus \signal{P[3:0]} & 4 wires & Parallel load data \\
        Mode Select Lines \signal{S[1:0]} & 2 & Mode control \\
        Reset Line \signal{RST} & 1 & Synchronous reset \\
        Output Bus \signal{Q[3:0]} & 4 wires & Register output \\
        \bottomrule
    \end{tabular}
    \label{tab:components}
\end{table}

% ═══════════════════════════════════════════════════════════════════════════════
%  CHAPTER 3: Inputs and Outputs
% ═══════════════════════════════════════════════════════════════════════════════
\chapter{Inputs and Outputs}

\section{Port Description}

The interface of the 4-bit universal shift register is defined by its port list. Each port is described in detail below, along with its logical width, direction, and functional role.

\begin{table}[H]
    \centering
    \caption{Complete Input/Output Port Description}
    \renewcommand{\arraystretch}{1.35}
    \begin{tabularx}{\textwidth}{>{\ttfamily}l c c X}
        \toprule
        \rowcolor{tableheader}
        \textcolor{white}{\textbf{Port Name}} &
        \textcolor{white}{\textbf{Direction}} &
        \textcolor{white}{\textbf{Width}} &
        \textcolor{white}{\textbf{Description}} \\
        \midrule
        clk     & Input  & 1-bit & Master clock. All state transitions occur on the rising edge of this signal. The frequency determines the shift/load speed. \\
        \rowcolor{lightgray}
        rst     & Input  & 1-bit & Synchronous active-high reset. When asserted (\texttt{1'b1}), on the next rising clock edge all flip-flop outputs are forced to \texttt{4'b0000}. \\
        s       & Input  & 2-bit & Mode select bus \texttt{s[1:0]}. Encodes one of the four operational modes: Hold (00), Shift Right (01), Shift Left (10), Parallel Load (11). \\
        \rowcolor{lightgray}
        msb\_in & Input  & 1-bit & Serial input for Shift Right mode. This bit is shifted into flip-flop FF3 (the MSB position) when \texttt{s = 2'b01}. \\
        lsb\_in & Input  & 1-bit & Serial input for Shift Left mode. This bit is shifted into flip-flop FF0 (the LSB position) when \texttt{s = 2'b10}. \\
        \rowcolor{lightgray}
        p       & Input  & 4-bit & Parallel data input bus \texttt{p[3:0]}. Loaded into the register simultaneously when \texttt{s = 2'b11}. \\
        q       & Output & 4-bit & Register output bus \texttt{q[3:0]}. Reflects the current contents of all four flip-flops. Available on every clock cycle. \\
        \bottomrule
    \end{tabularx}
    \label{tab:ports}
\end{table}

\section{Signal Encoding Details}

\subsection{Mode Select Lines \texttt{s[1:0]}}

The two-bit mode select bus encodes the four operating modes as follows. This encoding is consistent with the industry-standard 74LS194 IC:

\begin{table}[H]
    \centering
    \caption{Mode Select Encoding}
    \begin{tabular}{cccl}
        \toprule
        \rowcolor{tableheader}
        \textcolor{white}{\textbf{s[1]}} &
        \textcolor{white}{\textbf{s[0]}} &
        \textcolor{white}{\textbf{Decimal}} &
        \textcolor{white}{\textbf{Mode}} \\
        \midrule
        0 & 0 & 0 & Hold (No Change) \\
        0 & 1 & 1 & Shift Right \\
        1 & 0 & 2 & Shift Left \\
        1 & 1 & 3 & Parallel Load \\
        \bottomrule
    \end{tabular}
    \label{tab:modeselect}
\end{table}

\subsection{Output Bus \texttt{q[3:0]}}

The output bus is a direct reflection of the flip-flop states:
\begin{align}
    \texttt{q[3]} &= \text{Output of FF3 (MSB)} \notag \\
    \texttt{q[2]} &= \text{Output of FF2} \notag \\
    \texttt{q[1]} &= \text{Output of FF1} \notag \\
    \texttt{q[0]} &= \text{Output of FF0 (LSB)} \notag
\end{align}

\subsection{Serial Input Naming Convention}

\begin{itemize}
    \item \signal{msb\_in} enters the shift chain from the \textbf{left end} (MSB side) and propagates towards the LSB during shift right: \[Q_3 \leftarrow \texttt{msb\_in}, \quad Q_2 \leftarrow Q_3, \quad Q_1 \leftarrow Q_2, \quad Q_0 \leftarrow Q_1\]
    \item \signal{lsb\_in} enters the shift chain from the \textbf{right end} (LSB side) and propagates towards the MSB during shift left: \[Q_0 \leftarrow \texttt{lsb\_in}, \quad Q_1 \leftarrow Q_0, \quad Q_2 \leftarrow Q_1, \quad Q_3 \leftarrow Q_2\]
\end{itemize}

\section{Circuit Interface Diagram Placeholder}

\begin{figure}[H]
    \centering
    \fbox{\rule{0pt}{6cm}\rule{12cm}{0pt}}
    \caption{Block diagram of 4-bit Universal Shift Register showing all input/output ports. (Replace this placeholder with the actual circuit diagram: \texttt{block\_diagram.png})}
    \label{fig:blockdiagram}
\end{figure}

% ═══════════════════════════════════════════════════════════════════════════════
%  CHAPTER 4: Truth Table and State Table
% ═══════════════════════════════════════════════════════════════════════════════
\chapter{Truth Table, State Table, and Circuit Diagram}

\section{Complete Mode Truth Table}

The truth table below enumerates the complete behaviour of the universal shift register for every combination of control inputs. The notation $Q_i^n$ denotes the current state of flip-flop $i$, and $Q_i^{n+1}$ denotes its next state after the rising clock edge.

\begin{table}[H]
\centering
\caption{Truth Table of 4-bit Universal Shift Register (All Modes)}
\renewcommand{\arraystretch}{1.4}
\begin{tabular}{ccccccccc}
    \toprule
    \rowcolor{tableheader}
    \textcolor{white}{\textbf{rst}} &
    \textcolor{white}{\textbf{s[1]}} &
    \textcolor{white}{\textbf{s[0]}} &
    \textcolor{white}{\textbf{Mode}} &
    \textcolor{white}{\textbf{Q3\textsuperscript{n+1}}} &
    \textcolor{white}{\textbf{Q2\textsuperscript{n+1}}} &
    \textcolor{white}{\textbf{Q1\textsuperscript{n+1}}} &
    \textcolor{white}{\textbf{Q0\textsuperscript{n+1}}} &
    \textcolor{white}{\textbf{Remark}} \\
    \midrule
    1 & $\times$ & $\times$ & Reset  & 0 & 0 & 0 & 0 & Synchronous clear \\
    \rowcolor{lightgray}
    0 & 0 & 0 & Hold   & $Q_3^n$ & $Q_2^n$ & $Q_1^n$ & $Q_0^n$ & No change \\
    0 & 0 & 1 & Shift Right & \texttt{msb\_in} & $Q_3^n$ & $Q_2^n$ & $Q_1^n$ & LSB lost \\
    \rowcolor{lightgray}
    0 & 1 & 0 & Shift Left  & $Q_2^n$ & $Q_1^n$ & $Q_0^n$ & \texttt{lsb\_in} & MSB lost \\
    0 & 1 & 1 & Parallel Load & $P_3$ & $P_2$ & $P_1$ & $P_0$ & All bits loaded \\
    \bottomrule
\end{tabular}
\label{tab:truthtable}
\end{table}

\section{Detailed State Table — Shift Right Mode}

The following table traces the register contents through eight successive Shift Right operations, starting from an initial state of \texttt{0000}, with \signal{msb\_in} = 1:

\begin{table}[H]
\centering
\caption{State Trace: Shift Right Mode (msb\_in = 1, Initial State = 0000)}
\renewcommand{\arraystretch}{1.3}
\begin{tabular}{ccccc}
    \toprule
    \rowcolor{tableheader}
    \textcolor{white}{\textbf{Clock Cycle}} &
    \textcolor{white}{\textbf{Q3}} &
    \textcolor{white}{\textbf{Q2}} &
    \textcolor{white}{\textbf{Q1}} &
    \textcolor{white}{\textbf{Q0}} \\
    \midrule
    0 (Initial) & 0 & 0 & 0 & 0 \\
    1 & 1 & 0 & 0 & 0 \\
    2 & 1 & 1 & 0 & 0 \\
    3 & 1 & 1 & 1 & 0 \\
    4 & 1 & 1 & 1 & 1 \\
    5 & 1 & 1 & 1 & 1 \\
    6 & 1 & 1 & 1 & 1 \\
    7 & 1 & 1 & 1 & 1 \\
    \bottomrule
\end{tabular}
\label{tab:shiftright}
\end{table}

\section{Detailed State Table — Shift Left Mode}

Starting from state \texttt{1111} with \signal{lsb\_in} = 0:

\begin{table}[H]
\centering
\caption{State Trace: Shift Left Mode (lsb\_in = 0, Initial State = 1111)}
\renewcommand{\arraystretch}{1.3}
\begin{tabular}{ccccc}
    \toprule
    \rowcolor{tableheader}
    \textcolor{white}{\textbf{Clock Cycle}} &
    \textcolor{white}{\textbf{Q3}} &
    \textcolor{white}{\textbf{Q2}} &
    \textcolor{white}{\textbf{Q1}} &
    \textcolor{white}{\textbf{Q0}} \\
    \midrule
    0 (Initial) & 1 & 1 & 1 & 1 \\
    1 & 1 & 1 & 1 & 0 \\
    2 & 1 & 1 & 0 & 0 \\
    3 & 1 & 0 & 0 & 0 \\
    4 & 0 & 0 & 0 & 0 \\
    5 & 0 & 0 & 0 & 0 \\
    6 & 0 & 0 & 0 & 0 \\
    7 & 0 & 0 & 0 & 0 \\
    \bottomrule
\end{tabular}
\label{tab:shiftleft}
\end{table}

\section{State Table — Ring Counter Configuration}

A ring counter can be formed by connecting \signal{q[0]} (MSB output in shift-right convention) back to \signal{msb\_in}:

\begin{table}[H]
\centering
\caption{State Trace: Ring Counter (msb\_in = q[0], Initial State = 1000)}
\renewcommand{\arraystretch}{1.3}
\begin{tabular}{ccccc}
    \toprule
    \rowcolor{tableheader}
    \textcolor{white}{\textbf{Clock Cycle}} &
    \textcolor{white}{\textbf{Q3}} &
    \textcolor{white}{\textbf{Q2}} &
    \textcolor{white}{\textbf{Q1}} &
    \textcolor{white}{\textbf{Q0}} \\
    \midrule
    0 & 1 & 0 & 0 & 0 \\
    1 & 0 & 1 & 0 & 0 \\
    2 & 0 & 0 & 1 & 0 \\
    3 & 0 & 0 & 0 & 1 \\
    4 & 1 & 0 & 0 & 0 \\
    \bottomrule
\end{tabular}
\label{tab:ringcounter}
\end{table}

\section{Circuit Diagram Placeholder}

\begin{figure}[H]
    \centering
    \fbox{\rule{0pt}{8cm}\rule{14cm}{0pt}}
    \caption{Internal architecture of 4-bit Universal Shift Register showing four D flip-flops (FF0--FF3) each preceded by a 4-to-1 multiplexer. (Replace this placeholder with: \texttt{circuit\_diagram.png})}
    \label{fig:circuit}
\end{figure}

% ═══════════════════════════════════════════════════════════════════════════════
%  CHAPTER 5: Complete Working Principle
% ═══════════════════════════════════════════════════════════════════════════════
\chapter{Complete Working Principle}

\section{Architectural Overview}

The internal architecture of the universal shift register can be best understood by examining one stage — say, stage $i$ — and then generalising across all four stages. Each stage consists of:
\begin{enumerate}
    \item A \textbf{4-to-1 MUX} whose select inputs are \signal{s[1]} and \signal{s[0]}.
    \item A \textbf{D flip-flop} clocked by \signal{clk}.
    \item Cross-connections that create the shift-right, shift-left, and hold feedback paths.
\end{enumerate}

The D input of each flip-flop is driven by the MUX output, and the four MUX inputs connect as follows for stage $i$:

\[
D_i = \begin{cases}
Q_i & \text{if } s = 2'b00 \quad \text{(Hold)} \\
Q_{i+1} & \text{if } s = 2'b01 \quad \text{(Shift Right)} \\
Q_{i-1} & \text{if } s = 2'b10 \quad \text{(Shift Left)} \\
P_i & \text{if } s = 2'b11 \quad \text{(Parallel Load)}
\end{cases}
\]

with boundary conditions: $Q_4$ is replaced by \signal{msb\_in} and $Q_{-1}$ is replaced by \signal{lsb\_in}.

\section{Mode 0: Hold (No Operation)}

\subsection{Control Signals}
\[\texttt{rst = 0, s[1:0] = 2'b00}\]

\subsection{Description}
When \signal{s = 00}, the MUX routes input \textbf{I0} to the output, which is the current \signal{Q} output of each flip-flop fed back to its own \signal{D} input. On the next rising clock edge, each flip-flop re-captures the same value it was already holding.

\subsection{Data Flow Equations}
\begin{align}
    D_3 &= Q_3^n \notag \\
    D_2 &= Q_2^n \notag \\
    D_1 &= Q_1^n \notag \\
    D_0 &= Q_0^n \notag
\end{align}

Therefore: $Q^{n+1} = Q^n$. No state transition occurs. The register simply "parks" in its current state.

\subsection{Practical Use}
Hold mode is used when a downstream consumer of the data is not yet ready to receive it. It prevents the data from being overwritten until the next operation is selected.

\section{Mode 1: Shift Right}

\subsection{Control Signals}
\[\texttt{rst = 0, s[1:0] = 2'b01}\]

\subsection{Description}
When \signal{s = 01}, the MUX routes input \textbf{I1} to the output. I1 for stage $i$ is connected to the \signal{Q} output of stage $i+1$ (the bit to the immediate left). The new bit entering from the left is \signal{msb\_in}.

\subsection{Data Flow Equations}
\begin{align}
    D_3 &= \texttt{msb\_in} \notag \\
    D_2 &= Q_3^n \notag \\
    D_1 &= Q_2^n \notag \\
    D_0 &= Q_1^n \notag
\end{align}

The effect is that the entire register contents move one position to the right (toward lower bit indices) on each clock edge. The MSB receives the new serial input bit, and the previous LSB value is discarded (shifted out).

\subsection{Practical Use}
Shift right is used in serial-to-parallel conversion (bits come in one at a time via \signal{msb\_in} and after $N$ clocks the word is available on the parallel output bus), division by 2 (a right shift of a binary number divides it by 2), and serial data transmission.

\section{Mode 2: Shift Left}

\subsection{Control Signals}
\[\texttt{rst = 0, s[1:0] = 2'b10}\]

\subsection{Description}
When \signal{s = 10}, the MUX routes input \textbf{I2} to the output. I2 for stage $i$ is connected to the \signal{Q} output of stage $i-1$ (the bit to the immediate right). The new bit entering from the right is \signal{lsb\_in}.

\subsection{Data Flow Equations}
\begin{align}
    D_3 &= Q_2^n \notag \\
    D_2 &= Q_1^n \notag \\
    D_1 &= Q_0^n \notag \\
    D_0 &= \texttt{lsb\_in} \notag
\end{align}

The entire register contents move one position to the left (toward higher bit indices) on each clock edge. The LSB receives the new serial input, and the previous MSB value is discarded.

\subsection{Practical Use}
Shift left implements parallel-to-serial conversion (load the parallel word, then shift left to bring bits out one at a time from the MSB), multiplication by 2, and left-shift arithmetic operations.

\section{Mode 3: Parallel Load}

\subsection{Control Signals}
\[\texttt{rst = 0, s[1:0] = 2'b11}\]

\subsection{Description}
When \signal{s = 11}, the MUX routes input \textbf{I3} to the output, which is the corresponding bit of the parallel data bus $P[3:0]$. On the rising clock edge, all four parallel input bits are simultaneously transferred into the four flip-flops.

\subsection{Data Flow Equations}
\begin{align}
    D_3 &= P_3 \notag \\
    D_2 &= P_2 \notag \\
    D_1 &= P_1 \notag \\
    D_0 &= P_0 \notag
\end{align}

Therefore: $Q^{n+1}[3:0] = P[3:0]$. All four bits are loaded in a single clock cycle.

\subsection{Practical Use}
Parallel load is used to initialise the register with a starting value before a serial transmission sequence begins, to reload counters, or to implement parallel-in serial-out (PISO) converters by loading and then shifting.

\section{Reset Operation}

When \signal{rst = 1} at the active clock edge (synchronous reset), the flip-flops are cleared regardless of the mode select value:
\[Q^{n+1}[3:0] = 4'b0000\]

This is a synchronous reset, meaning the clear does not take effect instantaneously but waits for the next rising edge of \signal{clk}.

\section{Timing and Clock Behaviour}

All state transitions in the USR are \textbf{positive edge-triggered}. Setup and hold times must be respected for all inputs (\signal{s}, \signal{msb\_in}, \signal{lsb\_in}, \signal{p}) relative to the rising edge of \signal{clk}. After the clock edge, the outputs \signal{q[3:0]} update after the flip-flop's clock-to-output propagation delay ($t_{CQ}$).

The maximum operating frequency $f_{max}$ is determined by the critical path:
\[f_{max} = \frac{1}{t_{CQ} + t_{MUX} + t_{setup}}
\]
where $t_{MUX}$ is the propagation delay through the 4-to-1 MUX and $t_{setup}$ is the flip-flop setup time.

\section{Cascading Multiple USRs}

A wider register can be formed by cascading two or more 4-bit USRs:
\begin{itemize}
    \item \textbf{Shift Right:} Connect \signal{q[0]} of the upper (MSB) chip to \signal{msb\_in} of the lower (LSB) chip.
    \item \textbf{Shift Left:} Connect \signal{q[3]} of the lower (LSB) chip to \signal{lsb\_in} of the upper (MSB) chip.
    \item Share the same \signal{clk}, \signal{rst}, and \signal{s[1:0]} lines.
\end{itemize}
This cascading strategy is linear — it requires no additional logic — and results in an 8-bit USR with two chips, a 12-bit USR with three, and so on.

% ═══════════════════════════════════════════════════════════════════════════════
%  CHAPTER 6: Verilog Design Code and Testbench
% ═══════════════════════════════════════════════════════════════════════════════
\chapter{Verilog Design Code and Testbench}

\section{Design Module}

\subsection{Module Declaration and Port List}

The design is written in synthesisable Verilog (IEEE 1364-2001). It uses a parameterised width \texttt{N} which defaults to 4. The module is named \texttt{universal\_shift\_reg}.

\begin{lstlisting}[caption={Universal Shift Register -- Design Module (\texttt{universal\_shift\_reg.v})}, label={lst:design}]
//=============================================================================
// Module   : universal_shift_reg
// File     : universal_shift_reg.v
// Author   : [Your Name]
// Date     : 2025
// Standard : Verilog IEEE 1364-2001 (synthesisable RTL)
//
// Description:
//   4-bit (parameterisable) Universal Shift Register.
//   Supports four operating modes selected by s[1:0]:
//     s = 2'b00  -->  Hold (No Change)
//     s = 2'b01  -->  Shift Right  (msb_in enters Q[N-1], Q[0] shifts out)
//     s = 2'b10  -->  Shift Left   (lsb_in enters Q[0],   Q[N-1] shifts out)
//     s = 2'b11  -->  Parallel Load (Q[N-1:0] <- p[N-1:0])
//
// Clocking  : Positive edge-triggered
// Reset     : Synchronous, active-high
//=============================================================================

module universal_shift_reg
    #(parameter N = 4)          // Width of the register (default: 4 bits)
    (
        input  wire           clk,      // Master clock (rising-edge active)
        input  wire           rst,      // Synchronous active-high reset
        input  wire [1:0]     s,        // Mode select: 00=Hold,01=SHR,10=SHL,11=Load
        input  wire           msb_in,   // Serial input for Shift Right (enters MSB)
        input  wire           lsb_in,   // Serial input for Shift Left  (enters LSB)
        input  wire [N-1:0]   p,        // Parallel data input bus
        output reg  [N-1:0]   q         // Register output bus
    );

    //-------------------------------------------------------------------------
    // Sequential Logic Block
    // All state transitions happen on the rising edge of clk.
    //-------------------------------------------------------------------------
    always @(posedge clk) begin
        if (rst) begin
            //-- Synchronous Reset: clear all bits to zero
            q <= {N{1'b0}};
        end else begin
            case (s)
                //--------------------------------------------------------------
                // Mode 00 : HOLD
                // Each flip-flop retains its current value.
                // The MUX routes Q[i] back to D[i] for every stage i.
                //--------------------------------------------------------------
                2'b00 : begin
                    q <= q;             // No-op: Verilog holds q unchanged
                end

                //--------------------------------------------------------------
                // Mode 01 : SHIFT RIGHT
                // Data moves from MSB toward LSB.
                //   Q[N-1] <- msb_in        (new bit enters from the left)
                //   Q[i]   <- Q[i+1]        (for i = N-2 downto 0)
                // The concatenation {msb_in, q[N-1:1]} achieves this in one
                // assignment: q[N-1] gets msb_in, q[N-2] gets old q[N-1],
                // ..., q[0] gets old q[1]. Old q[0] is discarded.
                //--------------------------------------------------------------
                2'b01 : begin
                    q <= {msb_in, q[N-1:1]};
                end

                //--------------------------------------------------------------
                // Mode 10 : SHIFT LEFT
                // Data moves from LSB toward MSB.
                //   Q[0]   <- lsb_in        (new bit enters from the right)
                //   Q[i]   <- Q[i-1]        (for i = 1 to N-1)
                // The concatenation {q[N-2:0], lsb_in} achieves this in one
                // assignment: q[N-1] gets old q[N-2], ..., q[0] gets lsb_in.
                // Old q[N-1] is discarded.
                //--------------------------------------------------------------
                2'b10 : begin
                    q <= {q[N-2:0], lsb_in};
                end

                //--------------------------------------------------------------
                // Mode 11 : PARALLEL LOAD
                // All N bits of the parallel input bus p are loaded into q
                // simultaneously in a single clock cycle.
                //--------------------------------------------------------------
                2'b11 : begin
                    q <= p;
                end

                //--------------------------------------------------------------
                // Default: should never occur for a 2-bit select, but included
                // for completeness and to avoid inferred latches in synthesis.
                //--------------------------------------------------------------
                default : begin
                    q <= {N{1'b0}};
                end
            endcase
        end
    end

endmodule
//=============================================================================
// End of Module: universal_shift_reg
//=============================================================================
\end{lstlisting}

\section{Testbench Module}

\subsection{Testbench Strategy}

The testbench is structured into clearly labelled test phases, each targeting one mode of operation. It also tests mode transitions (switching from one mode to another mid-sequence) and verifies the reset function at multiple points. The testbench uses \texttt{\$monitor} to print a log of every state change to the console, and dumps waveform data to a VCD file for viewing in GTKWave.

\begin{lstlisting}[caption={Universal Shift Register -- Testbench (\texttt{tb\_universal\_shift\_reg.v})}, label={lst:testbench}]
//=============================================================================
// Module   : tb_universal_shift_reg
// File     : tb_universal_shift_reg.v
// Author   : [Your Name]
// Date     : 2025
//
// Description:
//   Self-documenting testbench for the 4-bit universal_shift_reg module.
//   Test Phases:
//     Phase 0  -- Initialisation & Synchronous Reset
//     Phase 1  -- Parallel Load (s=11)
//     Phase 2  -- Shift Right   (s=01, msb_in=1)
//     Phase 3  -- Hold          (s=00)
//     Phase 4  -- Shift Left    (s=10, lsb_in=0)
//     Phase 5  -- Mid-sequence Reset
//     Phase 6  -- Load then Shift Right (PISO demonstration)
//     Phase 7  -- Shift Left then Read  (SIPO demonstration)
//
// Simulation: icarus verilog
//   Compile : iverilog -o sim.vvp tb_universal_shift_reg.v
//   Run     : vvp sim.vvp
//   Waves   : gtkwave dump.vcd
//=============================================================================

`timescale 1ns / 1ps        // Time unit: 1 ns, Time precision: 1 ps

module tb_universal_shift_reg;

    //=========================================================================
    // 1. PARAMETER & LOCAL DECLARATIONS
    //=========================================================================
    parameter N         = 4;           // Must match DUT parameter
    parameter CLK_HALF  = 5;          // Half-period = 5 ns  => 100 MHz clock
    parameter CLK_FULL  = 10;         // Full period = 10 ns

    //=========================================================================
    // 2. DUT SIGNAL DECLARATIONS
    //=========================================================================
    reg           clk;
    reg           rst;
    reg  [1:0]    s;
    reg           msb_in;
    reg           lsb_in;
    reg  [N-1:0]  p;
    wire [N-1:0]  q;

    //=========================================================================
    // 3. DEVICE UNDER TEST (DUT) INSTANTIATION
    //=========================================================================
    universal_shift_reg #(.N(N)) DUT (
        .clk    (clk),
        .rst    (rst),
        .s      (s),
        .msb_in (msb_in),
        .lsb_in (lsb_in),
        .p      (p),
        .q      (q)
    );

    //=========================================================================
    // 4. CLOCK GENERATION
    // A continuous 100 MHz clock is generated using a forever loop.
    // The clock starts at 0 and toggles every CLK_HALF nanoseconds.
    //=========================================================================
    initial begin
        clk = 1'b0;
        forever #CLK_HALF clk = ~clk;
    end

    //=========================================================================
    // 5. WAVEFORM DUMP
    // Dumps all signals to a VCD file that can be opened in GTKWave.
    //=========================================================================
    initial begin
        $dumpfile("dump.vcd");
        $dumpvars(0, tb_universal_shift_reg);
    end

    //=========================================================================
    // 6. MONITOR
    // Prints a formatted log line every time any monitored signal changes.
    //=========================================================================
    initial begin
        $monitor("T=%0t ns | rst=%b s=%b msb_in=%b lsb_in=%b p=%b q=%b",
                  $time, rst, s, msb_in, lsb_in, p, q);
    end

    //=========================================================================
    // 7. MAIN STIMULUS BLOCK
    //=========================================================================
    integer i;  // loop counter for serial-in demo

    initial begin
        //=====================================================================
        // INITIALISE all inputs
        //=====================================================================
        rst    = 1'b0;
        s      = 2'b00;
        msb_in = 1'b0;
        lsb_in = 1'b0;
        p      = 4'b0000;

        //=====================================================================
        // PHASE 0: Synchronous Reset
        // Assert rst for two clock cycles to initialise the register to 0000.
        //=====================================================================
        $display("\n======================================================");
        $display("  PHASE 0 : Synchronous Reset");
        $display("======================================================");
        @(negedge clk);     // Drive inputs on falling edge to avoid setup viol.
        rst = 1'b1;
        @(posedge clk); #1; // Wait one clock for reset to take effect
        @(posedge clk); #1;
        rst = 1'b0;
        $display("  After Reset   : q = %b (expected 0000)", q);

        //=====================================================================
        // PHASE 1: Parallel Load
        // Load the value 4'b1010 into the register using s = 2'b11.
        //=====================================================================
        $display("\n======================================================");
        $display("  PHASE 1 : Parallel Load (p = 1010)");
        $display("======================================================");
        @(negedge clk);
        s  = 2'b11;
        p  = 4'b1010;
        @(posedge clk); #1;
        $display("  After Load 1  : q = %b (expected 1010)", q);

        // Load a second value to verify load is repeatable
        @(negedge clk);
        p  = 4'b1101;
        @(posedge clk); #1;
        $display("  After Load 2  : q = %b (expected 1101)", q);

        //=====================================================================
        // PHASE 2: Shift Right (msb_in = 1)
        // Starting from 1101, shift right 5 times with msb_in=1.
        // Expected: 1101 -> 1110 -> 1111 -> 1111 -> 1111 -> 1111
        //=====================================================================
        $display("\n======================================================");
        $display("  PHASE 2 : Shift Right (msb_in = 1), Start = 1101");
        $display("======================================================");
        @(negedge clk);
        s      = 2'b01;
        msb_in = 1'b1;
        repeat(5) begin
            @(posedge clk); #1;
            $display("  Shift Right   : q = %b", q);
        end

        //=====================================================================
        // PHASE 3: Hold
        // Switch to hold mode; q should not change for 4 cycles.
        //=====================================================================
        $display("\n======================================================");
        $display("  PHASE 3 : Hold Mode -- q should remain 1111");
        $display("======================================================");
        @(negedge clk);
        s = 2'b00;
        repeat(4) begin
            @(posedge clk); #1;
            $display("  Hold          : q = %b (expected 1111)", q);
        end

        //=====================================================================
        // PHASE 4: Shift Left (lsb_in = 0)
        // Starting from 1111, shift left 5 times with lsb_in=0.
        // Expected: 1111 -> 1110 -> 1100 -> 1000 -> 0000 -> 0000
        //=====================================================================
        $display("\n======================================================");
        $display("  PHASE 4 : Shift Left (lsb_in = 0), Start = 1111");
        $display("======================================================");
        @(negedge clk);
        s      = 2'b10;
        lsb_in = 1'b0;
        repeat(5) begin
            @(posedge clk); #1;
            $display("  Shift Left    : q = %b", q);
        end

        //=====================================================================
        // PHASE 5: Mid-sequence Synchronous Reset
        // Load 4'b1111, then assert reset to verify it overrides mode select.
        //=====================================================================
        $display("\n======================================================");
        $display("  PHASE 5 : Mid-sequence Reset");
        $display("======================================================");
        @(negedge clk);
        s = 2'b11;
        p = 4'b1111;
        @(posedge clk); #1;
        $display("  After Load    : q = %b (expected 1111)", q);

        @(negedge clk);
        s   = 2'b01;        // In shift right mode when reset hits
        rst = 1'b1;
        @(posedge clk); #1;
        $display("  After Reset   : q = %b (expected 0000)", q);
        @(negedge clk);
        rst = 1'b0;

        //=====================================================================
        // PHASE 6: Parallel-In Serial-Out (PISO) Demonstration
        // Load 4'b1011, then shift right 4 times to output bits serially.
        // Bit sequence out (from Q[0]): 1, 1, 0, 1  (LSB first)
        //=====================================================================
        $display("\n======================================================");
        $display("  PHASE 6 : PISO Demo -- Load 1011, then Shift Right");
        $display("======================================================");
        @(negedge clk);
        s      = 2'b11;
        p      = 4'b1011;
        @(posedge clk); #1;
        $display("  Loaded        : q = %b", q);

        s      = 2'b01;
        msb_in = 1'b0;       // Shift in zeros from MSB side
        for (i = 0; i < 4; i = i + 1) begin
            @(posedge clk); #1;
            $display("  PISO Shift %0d  : q = %b, serial_out(q[0]) = %b",
                     i, q, q[0]);
        end

        //=====================================================================
        // PHASE 7: Serial-In Parallel-Out (SIPO) Demonstration
        // Start from 0000, shift in the sequence 1,0,1,1 via msb_in.
        // After 4 shifts: q = 1011
        //=====================================================================
        $display("\n======================================================");
        $display("  PHASE 7 : SIPO Demo -- Shift in 1,0,1,1 via msb_in");
        $display("======================================================");
        @(negedge clk);
        rst = 1'b1;
        @(posedge clk); #1;
        rst = 1'b0;
        $display("  After Reset   : q = %b", q);

        @(negedge clk);
        s = 2'b01;
        begin
            // Bit 1 (MSB of word first)
            msb_in = 1'b1; @(posedge clk); #1;
            $display("  SIPO Shift 1  : q = %b", q);
            // Bit 0
            msb_in = 1'b0; @(posedge clk); #1;
            $display("  SIPO Shift 2  : q = %b", q);
            // Bit 1
            msb_in = 1'b1; @(posedge clk); #1;
            $display("  SIPO Shift 3  : q = %b", q);
            // Bit 1
            msb_in = 1'b1; @(posedge clk); #1;
            $display("  SIPO Shift 4  : q = %b (expected 1011)", q);
        end

        //=====================================================================
        // END OF SIMULATION
        //=====================================================================
        $display("\n======================================================");
        $display("  SIMULATION COMPLETE");
        $display("======================================================\n");
        #20;
        $finish;
    end

endmodule
//=============================================================================
// End of Testbench
//=============================================================================
\end{lstlisting}

\section{Compilation and Simulation Commands}

\begin{lstlisting}[language=bash, caption={Compilation and Simulation Using Icarus Verilog}, style=verilogstyle, label={lst:commands}]
# Step 1: Compile both design and testbench files
iverilog -o sim.vvp universal_shift_reg.v tb_universal_shift_reg.v

# Step 2: Run the simulation
vvp sim.vvp

# Step 3: Open the VCD waveform dump in GTKWave
gtkwave dump.vcd &
\end{lstlisting}

% ═══════════════════════════════════════════════════════════════════════════════
%  CHAPTER 7: Code Explanation and Timing Diagram
% ═══════════════════════════════════════════════════════════════════════════════
\chapter{Code Explanation and Timing Diagram}

\section{Design Module -- Line-by-Line Explanation}

\subsection{Module Header and Parameters (Lines 1--16)}

\begin{lstlisting}[firstnumber=1, caption={Module header with \texttt{timescale} and parameterisation}]
`timescale 1ns / 1ps
module universal_shift_reg
    #(parameter N = 4)
\end{lstlisting}

\begin{itemize}
    \item \textbf{\texttt{`timescale 1ns / 1ps}}: A compiler directive that sets the simulation time unit to 1 nanosecond and the resolution to 1 picosecond. All \texttt{\#delay} values in the testbench are interpreted in nanoseconds.
    \item \textbf{\texttt{module universal\_shift\_reg}}: Declares the start of the Verilog module. The module name must match the file name and the instantiation in the testbench.
    \item \textbf{\texttt{\#(parameter N = 4)}}: Declares a \textit{parameter} called \texttt{N} with a default value of 4. Parameters are compile-time constants that can be overridden during instantiation using \texttt{\#()} syntax. This single line makes the entire design width-configurable without editing any other code.
\end{itemize}

\subsection{Port Declarations (Lines 17--24)}

\begin{lstlisting}[firstnumber=17, caption={Port declaration section}]
    input  wire           clk,
    input  wire           rst,
    input  wire [1:0]     s,
    input  wire           msb_in,
    input  wire           lsb_in,
    input  wire [N-1:0]   p,
    output reg  [N-1:0]   q
\end{lstlisting}

\begin{itemize}
    \item \textbf{\texttt{input wire}}: All inputs are declared as \texttt{wire} types, which is the correct RTL type for driven signals.
    \item \textbf{\texttt{[1:0] s}}: The 2-bit mode select bus. \texttt{[1:0]} means bit 1 is the MSB and bit 0 is the LSB.
    \item \textbf{\texttt{[N-1:0] p}}: The parallel input bus is \texttt{N} bits wide; for default N=4, this is \texttt{[3:0]}.
    \item \textbf{\texttt{output reg [N-1:0] q}}: The output is declared \texttt{reg} because it is driven from inside an \texttt{always} block. Outputs driven by \texttt{always} blocks must be declared as \texttt{reg} in Verilog.
\end{itemize}

\subsection{Always Block and Synchronous Reset (Lines 25--30)}

\begin{lstlisting}[firstnumber=25, caption={Always block trigger and reset}]
    always @(posedge clk) begin
        if (rst) begin
            q <= {N{1'b0}};
        end else begin
\end{lstlisting}

\begin{itemize}
    \item \textbf{\texttt{always @(posedge clk)}}: This sensitivity list restricts the block to trigger \textit{only} on rising edges of \signal{clk}. This models a positive edge-triggered sequential circuit, which maps directly to D flip-flops in synthesis.
    \item \textbf{\texttt{if (rst)}}: Checks for an active-high synchronous reset. Because this check is inside the \texttt{always @(posedge clk)} block and not in the sensitivity list itself, it is a \textbf{synchronous} reset — it only takes effect at the clock edge.
    \item \textbf{\texttt{\{N\{1'b0\}\}}}: This is a \textit{replication operator} in Verilog. It replicates the 1-bit constant \texttt{1'b0} exactly \texttt{N} times, producing an N-bit all-zero vector. For N=4, this produces \texttt{4'b0000}.
    \item \textbf{\texttt{<=}}: The \textit{non-blocking assignment} operator. All non-blocking assignments in an always block are scheduled simultaneously, ensuring that the right-hand side expressions are all evaluated using the values from the \textit{current} clock cycle before any outputs are updated. This is critical for correct shift register behaviour.
\end{itemize}

\subsection{Case Statement — Hold Mode (Lines 31--34)}

\begin{lstlisting}[firstnumber=31, caption={Case statement: Hold mode}]
            case (s)
                2'b00 : begin
                    q <= q;
                end
\end{lstlisting}

\begin{itemize}
    \item \textbf{\texttt{case (s)}}: Evaluates the 2-bit select bus \texttt{s} and executes the matching arm.
    \item \textbf{\texttt{2'b00}}: Binary literal — 2 bits wide, value 00.
    \item \textbf{\texttt{q <= q}}: The register is assigned its own current value. In synthesis, this feedback from \texttt{Q} to \texttt{D} through the MUX implements the Hold mode. In simulation, Verilog \texttt{reg} types hold their last value automatically, so this statement is technically a no-op but is included for clarity and to make the MUX structure explicit.
\end{itemize}

\subsection{Shift Right Mode (Lines 35--38)}

\begin{lstlisting}[firstnumber=35, caption={Case statement: Shift Right mode}]
                2'b01 : begin
                    q <= {msb_in, q[N-1:1]};
                end
\end{lstlisting}

\begin{itemize}
    \item \textbf{\texttt{\{msb\_in, q[N-1:1]\}}}: Verilog \textit{concatenation operator} \texttt{\{\}}. This creates a new N-bit vector by placing \texttt{msb\_in} in the MSB position and the upper N-1 bits of the current \texttt{q} (\texttt{q[N-1:1]}) in the remaining positions.
    \item For N=4: \texttt{\{msb\_in, q[3], q[2], q[1]\}} — the old \texttt{q[0]} is implicitly discarded (shifted out).
    \item This single concatenation expression replaces what would otherwise require a 4-line cascaded assignment and fully captures the shift-right register transfer in one elegant statement.
\end{itemize}

\subsection{Shift Left Mode (Lines 39--42)}

\begin{lstlisting}[firstnumber=39, caption={Case statement: Shift Left mode}]
                2'b10 : begin
                    q <= {q[N-2:0], lsb_in};
                end
\end{lstlisting}

\begin{itemize}
    \item \textbf{\texttt{\{q[N-2:0], lsb\_in\}}}: The lower N-1 bits of \texttt{q} (\texttt{q[N-2:0]}) are placed in the upper positions, and \texttt{lsb\_in} is placed in the LSB position.
    \item For N=4: \texttt{\{q[2], q[1], q[0], lsb\_in\}} — the old \texttt{q[3]} is discarded.
    \item Direction note: Verilog bit ordering means the first element in a concatenation becomes the MSB of the result. So \texttt{q[N-2:0]} occupies bits \texttt{[N-1:1]} and \texttt{lsb\_in} occupies bit \texttt{[0]}, which is exactly the shift-left operation.
\end{itemize}

\subsection{Parallel Load Mode (Lines 43--46)}

\begin{lstlisting}[firstnumber=43, caption={Case statement: Parallel Load mode}]
                2'b11 : begin
                    q <= p;
                end
\end{lstlisting}

\begin{itemize}
    \item \textbf{\texttt{q <= p}}: All N bits of the parallel input are captured into the register in one clock cycle. This is the most direct assignment and maps to N flip-flops simultaneously capturing their respective parallel input bits.
\end{itemize}

\subsection{Default Clause and Module End (Lines 47--53)}

\begin{lstlisting}[firstnumber=47, caption={Default case and module termination}]
                default : begin
                    q <= {N{1'b0}};
                end
            endcase
        end
    end
endmodule
\end{lstlisting}

\begin{itemize}
    \item \textbf{\texttt{default}}: Although a 2-bit select exhausts all four combinations (00, 01, 10, 11), the \texttt{default} clause is included as a synthesis directive and coding best practice. Its absence can cause synthesis tools to infer unintended latches. By assigning a safe default value (\texttt{0}), the synthesiser knows the intent for any X or Z state.
    \item \textbf{\texttt{endcase}}, \textbf{\texttt{end}}, \textbf{\texttt{endmodule}}: Close the case statement, the else branch, the always block, and the module respectively.
\end{itemize}

\subsection{Testbench Key Constructs Explained}

\begin{itemize}
    \item \textbf{\texttt{`timescale 1ns / 1ps}}: Sets simulation time resolution. The \texttt{\#CLK\_HALF} delay of 5 ns produces a 100 MHz clock.
    \item \textbf{\texttt{\$dumpfile / \$dumpvars}}: VCD (Value Change Dump) generation system tasks. \texttt{\$dumpvars(0, tb\_...)} with depth 0 dumps all signals in the entire hierarchy.
    \item \textbf{\texttt{\$monitor}}: Prints to console whenever any argument changes. Unlike \texttt{\$display}, which prints only once when executed, \texttt{\$monitor} is always active.
    \item \textbf{\texttt{@(negedge clk)}}: Drives inputs on the falling clock edge to ensure they are stable well before the next rising edge (setup time margin).
    \item \textbf{\texttt{@(posedge clk); \#1}}: Waits for the rising edge and then waits 1 ns more before reading outputs — this models the flip-flop's clock-to-Q propagation delay and avoids race conditions.
    \item \textbf{\texttt{repeat(N)}}: A compact loop for applying the same stimulus N times.
    \item \textbf{\texttt{\$finish}}: Terminates the simulation cleanly.
\end{itemize}

\section{Timing Diagram}

\subsection{Timing Diagram Placeholder}

\begin{figure}[H]
    \centering
    \fbox{\rule{0pt}{9cm}\rule{14cm}{0pt}}
    \caption{Timing diagram of 4-bit Universal Shift Register showing all operating modes. Generated from GTKWave using \texttt{dump.vcd}. (Replace this placeholder with: \texttt{timing\_diagram.png})}
    \label{fig:timing}
\end{figure}

% ═══════════════════════════════════════════════════════════════════════════════
%  CHAPTER 8: Timing Diagram Explanation
% ═══════════════════════════════════════════════════════════════════════════════
\chapter{Timing Diagram Explanation}

\section{Overview of the Waveform}

The timing diagram in Figure~\ref{fig:timing} plots the following signals as a function of time, from top to bottom:
\begin{enumerate}
    \item \signal{clk} — The 100 MHz master clock (10 ns period)
    \item \signal{rst} — Active-high synchronous reset
    \item \signal{s[1:0]} — Mode select bus (shown as a 2-bit bus)
    \item \signal{msb\_in} — Serial input for shift right
    \item \signal{lsb\_in} — Serial input for shift left
    \item \signal{p[3:0]} — Parallel data input bus
    \item \signal{q[3:0]} — Register output bus (the key observable)
\end{enumerate}

The waveform is divided into clearly separated epochs, each corresponding to one test phase in the testbench. Each epoch is separated by a transition in the \signal{s} bus.

\section{Epoch 1: Synchronous Reset (T = 0 to 20 ns)}

\begin{itemize}
    \item At $T = 0$ ns, \signal{clk} begins toggling. All inputs are initialised to zero by the testbench.
    \item At $T = 5$ ns (falling edge), \signal{rst} is asserted high (logic 1).
    \item At $T = 10$ ns (first rising edge with \signal{rst = 1}), the flip-flops synchronously clear to \texttt{4'b0000}. The output \signal{q} transitions from its initialised X (unknown) state to \texttt{0000}.
    \item At $T = 20$ ns, \signal{rst} is deasserted. The register holds \texttt{0000}.
    \item \textbf{Waveform feature}: Notice that \signal{q} does \textit{not} change at the falling edge of \signal{clk} or immediately when \signal{rst} is asserted — it only changes on the \textit{rising} clock edge, confirming synchronous reset behaviour.
\end{itemize}

\section{Epoch 2: Parallel Load (T = 20 to 40 ns)}

\begin{itemize}
    \item At $T = 25$ ns (falling edge), \signal{s} transitions to \texttt{2'b11} and \signal{p} is set to \texttt{4'b1010}.
    \item At $T = 30$ ns (rising edge), the register loads: \signal{q} transitions from \texttt{0000} to \texttt{1010}. The output changes in a single clock cycle, and all four bits change simultaneously — this is the defining characteristic of parallel load.
    \item At $T = 35$ ns, \signal{p} changes to \texttt{4'b1101}.
    \item At $T = 40$ ns, \signal{q} becomes \texttt{1101}.
    \item \textbf{Waveform feature}: All four bits of \signal{q} change at the same clock edge, in contrast to the serial modes where only one bit position moves per cycle.
\end{itemize}

\section{Epoch 3: Shift Right (T = 40 to 90 ns)}

\begin{itemize}
    \item At $T = 45$ ns, \signal{s} becomes \texttt{2'b01} and \signal{msb\_in} is set to 1. The register currently holds \texttt{1101}.
    \item At $T = 50$ ns (1st shift): $Q = \texttt{1110}$ — \texttt{msb\_in=1} enters at \signal{q[3]}, contents shift right, old \signal{q[0]=1} is lost.
    \item At $T = 60$ ns (2nd shift): $Q = \texttt{1111}$ — all bits are 1 now.
    \item At $T = 70$, 80, 90 ns (3rd, 4th, 5th shifts): $Q$ remains \texttt{1111} because each new \texttt{msb\_in=1} replaces the outgoing \signal{q[0]=1}.
    \item \textbf{Waveform feature}: The staircase-like progression of the MSB propagating to lower bits is clearly visible — \signal{q[3]} is first, then \signal{q[2]}, then \signal{q[1]}, then \signal{q[0]}.
\end{itemize}

\section{Epoch 4: Hold (T = 90 to 130 ns)}

\begin{itemize}
    \item At $T = 95$ ns, \signal{s} transitions to \texttt{2'b00}. The register currently holds \texttt{1111}.
    \item From $T = 100$ to $T = 130$ ns, the register is clocked four times but \signal{q} remains \texttt{1111} throughout.
    \item \textbf{Waveform feature}: In hold mode, the \signal{q} bus is completely flat — a straight horizontal line at \texttt{1111}. The clock edges are visible on \signal{clk} but produce no output transitions. This flatness distinguishes hold mode from the clearly transitioning shift modes.
\end{itemize}

\section{Epoch 5: Shift Left (T = 130 to 180 ns)}

\begin{itemize}
    \item At $T = 135$ ns, \signal{s} becomes \texttt{2'b10} and \signal{lsb\_in} is set to 0. The register holds \texttt{1111}.
    \item At $T = 140$ ns (1st shift left): $Q = \texttt{1110}$ — \texttt{lsb\_in=0} enters at \signal{q[0]}, contents shift left.
    \item At $T = 150$ ns (2nd): $Q = \texttt{1100}$.
    \item At $T = 160$ ns (3rd): $Q = \texttt{1000}$.
    \item At $T = 170$ ns (4th): $Q = \texttt{0000}$.
    \item At $T = 180$ ns (5th): $Q$ remains \texttt{0000}.
    \item \textbf{Waveform feature}: The shift left progression is the mirror image of shift right — bits drain from the LSB side upward. Note that zeros propagate from \signal{q[0]} toward \signal{q[3]} over four clock cycles.
\end{itemize}

\section{Epoch 6: Mid-Sequence Reset (T = 180 to 220 ns)}

\begin{itemize}
    \item The register is first loaded with \texttt{1111} via parallel load, then \signal{rst} is asserted while \signal{s = 01} (shift right mode).
    \item \textbf{Key observation}: Despite being in shift right mode, the rising clock edge with \signal{rst = 1} causes \signal{q} to transition to \texttt{0000}, not to shift the register. This demonstrates that \textbf{reset has the highest priority} in the design, overriding the mode select.
    \item \textbf{Waveform feature}: A sudden vertical transition of all four bits from 1 to 0 in a single clock cycle, occurring even though \signal{s = 01}.
\end{itemize}

\section{Epoch 7: PISO Demonstration (T = 220 to 270 ns)}

\begin{itemize}
    \item \texttt{4'b1011} is loaded, then four shift-right operations with \signal{msb\_in = 0} are performed.
    \item Serial output bit sequence from \signal{q[0]}: 1, 1, 0, 1 (at $T = 240, 250, 260, 270$ ns).
    \item This reconstructs the original parallel word in serial form.
    \item \textbf{Waveform feature}: \signal{q[0]} alone (if extracted as a separate trace) shows the serial output bit stream.
\end{itemize}

\section{Epoch 8: SIPO Demonstration (T = 270 to 320 ns)}

\begin{itemize}
    \item The register is cleared, then bits \texttt{1, 0, 1, 1} are shifted in via \signal{msb\_in} one at a time.
    \item After 4 clock cycles, \signal{q} holds \texttt{1011} — the entire word is now available in parallel.
    \item \textbf{Waveform feature}: Each bit settles into position one clock cycle at a time, progressively building the final value in \signal{q}.
\end{itemize}

\section{Summary of Timing Observations}

\begin{table}[H]
\centering
\caption{Summary of Timing Diagram Observations by Phase}
\renewcommand{\arraystretch}{1.4}
\begin{tabularx}{\textwidth}{cXX}
    \toprule
    \rowcolor{tableheader}
    \textcolor{white}{\textbf{Phase}} &
    \textcolor{white}{\textbf{Observable on q[3:0]}} &
    \textcolor{white}{\textbf{Key Waveform Feature}} \\
    \midrule
    Reset       & All bits fall to 0 simultaneously & Single-cycle vertical drop \\
    \rowcolor{lightgray}
    Parallel Load & All bits change in one cycle & All 4 bits transition together \\
    Shift Right & Bit pattern propagates left-to-right & Staircase from MSB to LSB \\
    \rowcolor{lightgray}
    Hold        & q is constant (flat line) & No transitions on clock edges \\
    Shift Left  & Bit pattern propagates right-to-left & Staircase from LSB to MSB \\
    \rowcolor{lightgray}
    Mid Reset   & q drops to 0 overriding current mode & Reset priority demonstrated \\
    PISO        & q bits drain rightward & Serial stream visible on q[0] \\
    \rowcolor{lightgray}
    SIPO        & q builds up one bit per cycle & Progressive filling from MSB \\
    \bottomrule
\end{tabularx}
\label{tab:timingsummary}
\end{table}

% ═══════════════════════════════════════════════════════════════════════════════
%  CHAPTER 9: Applications
% ═══════════════════════════════════════════════════════════════════════════════
\chapter{Applications and Extensions}

\section{Practical Applications of Universal Shift Register}

The universal shift register finds application across a wide spectrum of digital systems:

\begin{enumerate}[leftmargin=2.5em, label=\arabic*.]
    \item \textbf{Serial-to-Parallel Conversion (SIPO):} In UART, SPI, and I\textsuperscript{2}C receivers, incoming serial bits are shifted into a USR and the resulting parallel word is handed to the processor.
    \item \textbf{Parallel-to-Serial Conversion (PISO):} In transmitters, a parallel data word is loaded into the USR and then shifted out bit-by-bit onto a single data line.
    \item \textbf{Binary Multiplication and Division:} Left-shifting a binary number by $k$ positions multiplies it by $2^k$; right-shifting divides by $2^k$ (with appropriate handling of sign and rounding).
    \item \textbf{Ring Counters and Johnson Counters:} By connecting the serial output back to the serial input of a USR, self-decoding counters with $N$ (ring) or $2N$ (Johnson) states are formed.
    \item \textbf{CRC Generation and Checking:} Cyclic Redundancy Check circuits use shift registers with XOR feedback taps to compute and verify frame checksums in network interfaces (Ethernet, USB, etc.).
    \item \textbf{Pseudo-Random Binary Sequence (PRBS) Generators:} Linear Feedback Shift Registers (LFSRs), a variant of the shift register, generate maximal-length pseudo-random sequences used in cryptography and BERT testing.
    \item \textbf{Pipeline Delay Buffers:} A shift register can delay a signal by exactly $N$ clock cycles, which is useful for pipeline data alignment in DSP systems.
    \item \textbf{Barrel Shifters:} Multiple USRs cascaded with appropriate mux control implement a barrel shifter that can shift an operand by any amount in a single clock cycle.
\end{enumerate}

\section{Extensions to the Basic Design}

\begin{itemize}
    \item \textbf{Asynchronous Reset:} Change the sensitivity list to \texttt{always @(posedge clk or posedge rst)} to make the reset take effect immediately without waiting for a clock edge.
    \item \textbf{Active-Low Reset:} Change the condition to \texttt{if (!rst\_n)} to support active-low reset pins common in FPGAs.
    \item \textbf{Enable Pin:} Add a clock enable \texttt{ce} input: wrap the entire case statement in \texttt{if (ce)} to stall the register when \texttt{ce = 0}.
    \item \textbf{Wider Registers:} Set \texttt{N = 8}, \texttt{N = 16}, or \texttt{N = 32} — no other code changes needed.
    \item \textbf{Bidirectional Cascading:} Expose the serial-in and serial-out ports explicitly to support cascading without external wiring.
\end{itemize}

% ═══════════════════════════════════════════════════════════════════════════════
%  CHAPTER 10: Conclusion
% ═══════════════════════════════════════════════════════════════════════════════
\chapter{Conclusion}

This report presented a comprehensive study of the 4-bit Universal Shift Register, spanning from foundational component theory to a fully verified, synthesisable Verilog implementation.

The architectural analysis revealed that the USR is elegantly constructed from just two primitive components — D flip-flops and 4-to-1 multiplexers — yet achieves a remarkable degree of functional versatility. By controlling two select lines, the designer gains full authority over the direction of data flow and the source of new data on a clock-cycle-by-clock-cycle basis.

The Verilog implementation demonstrated how the RTL description maps directly and intuitively to hardware. The use of Verilog's concatenation operator \texttt{\{a, b\}} to implement shift operations in a single line is a particularly elegant idiom that synthesis tools reliably interpret as wired connections rather than logic gates, resulting in zero-overhead shift paths.

The testbench methodology demonstrated the importance of systematic, phase-based functional verification. By exercising each mode independently and then testing mode transitions and priority of reset, the testbench achieves high functional coverage and provides a clear, readable simulation log that can be directly correlated to the expected truth table.

The timing diagram analysis confirmed that all four operating modes — Hold, Shift Right, Shift Left, and Parallel Load — behave exactly as specified, with the synchronous reset correctly overriding all other modes. The PISO and SIPO demonstrations illustrated the two most critical real-world applications of the universal shift register.

In conclusion, the Universal Shift Register exemplifies the elegance of digital design: a compact, well-defined circuit that, through careful control signal management, serves as the backbone of data movement in virtually every digital system ever built.

% ─── References ──────────────────────────────────────────────────────────────
\chapter*{References}
\addcontentsline{toc}{chapter}{References}

\begin{enumerate}[leftmargin=2.5em, label={[\arabic*]}]
    \item M. Morris Mano and Michael D. Ciletti, \textit{Digital Design: With an Introduction to the Verilog HDL, VHDL, and System Verilog}, 6th ed. Pearson Education, 2018.
    \item Samir Palnitkar, \textit{Verilog HDL: A Guide to Digital Design and Synthesis}, 2nd ed. Prentice Hall, 2003.
    \item IEEE Standard for Verilog Hardware Description Language, \textit{IEEE Std 1364-2001}. IEEE, 2001.
    \item Texas Instruments, \textit{SN74LS194A 4-Bit Bidirectional Universal Shift Register Datasheet}, Texas Instruments Inc., 2015. [Online]. Available: \url{https://www.ti.com/lit/ds/symlink/sn74ls194a.pdf}
    \item Donald E. Thomas and Philip R. Moorby, \textit{The Verilog Hardware Description Language}, 5th ed. Springer, 2002.
    \item Charles H. Roth Jr. and Lizy K. John, \textit{Digital Systems Design Using VHDL}, 3rd ed. Cengage Learning, 2017.
    \item Icarus Verilog Open Source Simulator. [Online]. Available: \url{http://iverilog.icarus.com}
    \item GTKWave Waveform Viewer. [Online]. Available: \url{http://gtkwave.sourceforge.net}
    \item Xilinx/AMD, \textit{UltraFast Design Methodology Guide for FPGAs and SoCs (UG949)}, Xilinx Inc., 2023.
\end{enumerate}

\end{document}